\section{How to Work with This Course}

This book is part of a weekly learning process rather than a collection of
material to be read only before an examination. Each of the twelve main
chapters is designed to support one week of study. The chapter provides the
mathematical structure for the lecture, a complete set of notes to which the
student can return, and a starting point for individual problem solving and
classroom discussion.

The purpose of this organisation is to make mathematical work continuous and
visible. A student should not leave a topic behind when the lecture ends. The
lecture introduces the main questions and guides the first encounter with the
ideas. The written chapter then makes it possible to revisit the reasoning at a
slower pace, check definitions and examples, and identify points that require
further thought.

\subsection{The Weekly Learning Cycle}

A typical week in this course consists of the following connected stages.

\begin{enumerate}
    \item \textbf{Attend the lecture.}
    The lecture introduces the topic, motivates the central questions, and
    demonstrates how the relevant mathematical objects and methods should be
    read.

    \item \textbf{Return to the chapter.}
    After the lecture, read the corresponding chapter carefully. The PDF is
    not merely a summary of slides. It is an organised narrative containing
    definitions, explanations, examples, interpretations, and mathematical
    arguments that can be examined more than once.

    \item \textbf{Ask questions and clarify the material.}
    Mark every place where a definition, transition, calculation, or conclusion
    is not yet clear. Use the text, your own examples, classroom discussion,
    and appropriate digital tools to resolve these questions.

    \item \textbf{Solve the accompanying problems.}
    Use the problems to move from following an explanation to constructing an
    argument independently. A complete solution should state what is known,
    identify the method being used, justify the important steps, and interpret
    the result when interpretation is meaningful.

    \item \textbf{Prepare professional mathematical notes.}
    Organise your work so that another person can read it. Use precise notation,
    complete sentences where explanation is needed, and a clear sequence of
    reasoning. The final notes should show not only the answer, but also how the
    answer follows from the definitions and earlier results.

    \item \textbf{Present and discuss your work.}
    During exercise classes, be prepared to explain your solution, answer
    questions about individual steps, and describe what you learned while
    working on the problem. Presenting mathematics is part of learning
    mathematics: it reveals whether an argument has genuinely been understood.
\end{enumerate}

These stages should not be treated as independent obligations. They form one
cycle:
\[
\text{lecture}
\longrightarrow
\text{careful reading}
\longrightarrow
\text{questions}
\longrightarrow
\text{problem solving}
\longrightarrow
\text{written explanation}
\longrightarrow
\text{presentation}.
\]
Each stage prepares the next one, and the classroom presentation provides
feedback for further reading and revision.

\subsection{Learning from Feedback to the Group}

Feedback given during exercise classes is intended for the whole group, not
only for the student whose work is currently being presented. When another
student explains a solution, everyone should observe the presentation, follow
the argument, and listen carefully to the feedback that follows.

By comparing different examples, students can learn what a complete solution
looks like, which explanations make an argument clear, which steps require
justification, and which weaknesses cause a note to be returned for revision.
Feedback on one piece of work may reveal a principle that applies to many other
solutions.

Students should therefore not focus only on comments addressed directly to
their own work. Listening to the presentations and feedback of others provides
a broader set of examples from which to understand the standards of the course.
It helps the group develop an operational understanding of what is expected:
what makes mathematical work clear, justified, and complete, and what still
needs improvement.

The feedback concerns the mathematical work and its results, not the personal
worth of the student. What is examined is the quality of the note, the
understanding demonstrated in the argument, and the clarity of the
presentation. Feedback is intended to make expectations visible, identify what
already works, and show what should be changed in order to produce a better
solution.

The first submitted version does not always have to be the final version.
Students should learn to receive comments, reconsider their choices, correct
errors, add missing explanations, and improve both the written note and the way
in which it is presented. Revision is not evidence of failure. It is a normal
part of serious intellectual work.

This creates a recurring process:
\[
\text{problem}
\longrightarrow
\text{solution}
\longrightarrow
\text{presentation}
\longrightarrow
\text{feedback}
\longrightarrow
\text{revision}.
\]
The result should be a stronger piece of work: a clearer note, deeper
understanding, and a more precise presentation.

This process is important throughout university study and later in
professional life. Students and professionals are given tasks, asked to solve
problems, expected to present their results, and required to respond to the
opinions and criticism of others. Learning to improve an initial solution
through feedback is therefore not an additional classroom exercise. It is a
fundamental habit of responsible academic and professional work.

\subsection{The Role of AI}

Artificial intelligence may be used as an individual learning assistant. It can
help explain a difficult paragraph, generate a simpler example, ask diagnostic
questions, examine a proposed solution, point out missing justifications, and
improve the organisation of mathematical writing. It can also help a student
turn rough calculations into a clear and professional solution.

AI does not, however, replace the student's mathematical work. Receiving a
well-written answer is not the same as understanding it. The student remains
responsible for checking definitions, calculations, assumptions, and
conclusions. Any AI-assisted solution included in the student's notes must be
understood well enough to be reconstructed, explained, and defended in class
without relying on the tool to speak on the student's behalf.

A productive conversation with AI should therefore contain questions such as:
Why is this step valid? Which definition is being used? What would change if an
assumption were removed? Can the same result be explained geometrically? Is
there a counterexample to the statement I have written? Such questions use AI
to deepen reasoning rather than to avoid it.

\subsection{What Progress Looks Like}

Progress in this course is demonstrated by the quality of the student's
mathematical work. By the end of a weekly cycle, the student should be able to
identify the main objects introduced in the chapter, state the essential
definitions, solve representative problems, justify the principal steps, and
communicate the resulting argument clearly.

The twelve thematic chapters provide the main rhythm of the semester. The
remaining course time can be used to introduce the method of work, connect
topics, revise important ideas, and present more substantial solutions. The
goal is that the student finishes the course not only with a collection of
results, but also with a set of carefully prepared notes and the ability to
read, write, question, and present mathematics with increasing independence.
